% GENERATED FILE. Edit canonical lessons, then run npm run generate:books.
% canonical-source-hash: fnv1a64:d0d67bb05bd42fbe
% canonical-lessons: PA-C290-asman, PA-C290-gare, PA-C290-khullhna, PA-C290-ghatna, PA-C290-vagna

\chapter{The Week's Weather}
\label{ch:pa-a2-the-week-s-weather}
\hlchaptermodality{290}

\begin{chapteropening}
\textbf{By the end of this chapter:} \emph{I can say how the weather will change this week.}

Complete the last lesson of chapter 290: I can say how the weather will change this week.
\end{chapteropening}

\section[\texorpdfstring{\pa{ਅਸਮਾਨ} (asmān)}{ਅਸਮਾਨ (asmān)}]{\pa{ਅਸਮਾਨ} (asmān) — the sky}

\label{lesson:PA-C290-asman}



\begin{quote}
Before the new one: say the Punjabi for to be chilled, then the Punjabi for the cold season.
\end{quote}

\subsection*{You'll want to know: \pa{ਅਸਮਾਨ}}
\textbf{\pa{ਅਸਮਾਨ}} — \emph{asmān} — \textquotedblleft{}the sky\textquotedblright{}.

\begin{quote}
\pa{ਅੱਜ} \pa{ਅਸਮਾਨ} \pa{ਨੀਲਾ} \pa{ਹੈ।}
\end{quote}

\begin{quote}
\emph{ajj asmān nīlā hai.}
\end{quote}

\begin{quote}
Today the sky is blue.
\end{quote}

A masculine noun.

\begin{cousinweb}
\textbf{\pa{ਅਸਮਾਨ}} (\emph{asmān}) is a loan, usually traced to Persian \emph{āsmān}, \textquotedblleft{}sky\textquotedblright{}.
\end{cousinweb}

\begin{grammarlens}[title={What you've built}]
One more word for saying what comes next.
\end{grammarlens}

\subsection*{Your turn}
\begin{itemize}
\raggedright
  \item \emph{Say it:} \emph{asmān}
  \item \emph{Say it:} \emph{asmān}, once more
  \item \emph{Say it:} \emph{ajj asmān nīlā hai.} — Today the sky is blue.
  \item \emph{Recall:} say the Punjabi for to be chilled, then the Punjabi for the cold season, then say \emph{asmān} again
\end{itemize}

\begin{tcolorbox}[breakable,colback=teal!4,colframe=teal!35!black,title={Before you move on}]
What does \pa{ਅਸਮਾਨ} mean? (\textquotedblleft{}The sky\textquotedblright{}.) Say it once more.
\end{tcolorbox}
\section[\texorpdfstring{\pa{ਗੜੇ} (gaṛe)}{ਗੜੇ (gaṛe)}]{\pa{ਗੜੇ} (gaṛe) — hail; hailstones}

\label{lesson:PA-C290-gare}



\begin{quote}
Before the new one: say the Punjabi for the cold season, then the Punjabi for the sky.
\end{quote}

\subsection*{You'll want to know: \pa{ਗੜੇ}}
\textbf{\pa{ਗੜੇ}} — \emph{gaṛe} — \textquotedblleft{}hail; hailstones\textquotedblright{}.

\begin{quote}
\pa{ਕੱਲ੍ਹ} \pa{ਗੜੇ} \pa{ਪੈਣਗੇ।}
\end{quote}

\begin{quote}
\emph{kall gaṛe paiṇge.}
\end{quote}

\begin{quote}
There will be hail tomorrow.
\end{quote}

A masculine noun, used in the plural.

\begin{cousinweb}
\textbf{\pa{ਗੜੇ}} (\emph{gaṛe}) is inherited from a reconstructed \emph{*gaḍa}, \textquotedblleft{}dropping\textquotedblright{}; CDIAL lists Punjabi \emph{gaṛā}, \textquotedblleft{}hail\textquotedblright{}, under it (CDIAL 3969).
\end{cousinweb}

\begin{grammarlens}[title={What you've built}]
One more word for saying what comes next.
\end{grammarlens}

\subsection*{Your turn}
\begin{itemize}
\raggedright
  \item \emph{Say it:} \emph{gaṛe}
  \item \emph{Say it:} \emph{gaṛe}, once more
  \item \emph{Say it:} \emph{kall gaṛe paiṇge.} — There will be hail tomorrow.
  \item \emph{Recall:} say the Punjabi for the cold season, then the Punjabi for the sky, then say \emph{gaṛe} again
  \item \emph{Write it:} copy \textbf{\pa{ਗੜੇ}} once, checking the whole word
\end{itemize}

\begin{tcolorbox}[breakable,colback=teal!4,colframe=teal!35!black,title={Before you move on}]
What does \pa{ਗੜੇ} mean? (\textquotedblleft{}Hail; hailstones\textquotedblright{}.) Say it once more.
\end{tcolorbox}
\section[\texorpdfstring{\pa{ਖੁੱਲ੍ਹਣਾ} (khullhṇā)}{ਖੁੱਲ੍ਹਣਾ (khullhṇā)}]{\pa{ਖੁੱਲ੍ਹਣਾ} (khullhṇā) — to clear up; to come open}

\label{lesson:PA-C290-khullhna}



\begin{quote}
Before the new one: say the Punjabi for the sky, then the Punjabi for hail.
\end{quote}

\subsection*{You'll want to know: \pa{ਖੁੱਲ੍ਹਣਾ}}
\textbf{\pa{ਖੁੱਲ੍ਹਣਾ}} — \emph{khullhṇā} — \textquotedblleft{}to clear up; to come open\textquotedblright{}.

\begin{quote}
\pa{ਦੁਪਹਿਰ} \pa{ਤੱਕ} \pa{ਮੌਸਮ} \pa{ਖੁੱਲ੍ਹੇਗਾ।}
\end{quote}

\begin{quote}
\emph{dupahir tak mausam khullhegā.}
\end{quote}

\begin{quote}
The weather will clear up by noon.
\end{quote}

Past (it …): \pa{ਉਹ} \pa{ਦਰਵਾਜ਼ਾ} \pa{ਖੁੱਲ੍ਹਿਆ}, \emph{oh darvāzā khullhiā}.

\begin{cousinweb}
\textbf{\pa{ਖੁੱਲ੍ਹਣਾ}} (\emph{khullhṇā}) is inherited from a reconstructed \emph{*khull}, \textquotedblleft{}to be open\textquotedblright{}; CDIAL lists Punjabi \emph{khullhṇā}, \textquotedblleft{}to be open\textquotedblright{}, under it (CDIAL 3945).
\end{cousinweb}

\begin{grammarlens}[title={What you've built}]
One more word for saying what comes next.
\end{grammarlens}

\subsection*{Your turn}
\begin{itemize}
\raggedright
  \item \emph{Say it:} \emph{khullhṇā}
  \item \emph{Say it:} \emph{khullhṇā}, once more
  \item \emph{Say it:} \emph{dupahir tak mausam khullhegā.} — The weather will clear up by noon.
  \item \emph{Recall:} say the Punjabi for the sky, then the Punjabi for hail, then say \emph{khullhṇā} again
  \item \emph{Write it:} copy \textbf{\pa{ਖੁੱਲ੍ਹਣਾ}} once, checking the whole word
\end{itemize}

\begin{tcolorbox}[breakable,colback=teal!4,colframe=teal!35!black,title={Before you move on}]
What does \pa{ਖੁੱਲ੍ਹਣਾ} mean? (\textquotedblleft{}To clear up; to come open\textquotedblright{}.) Say it once more.
\end{tcolorbox}
\section[\texorpdfstring{\pa{ਘਟਣਾ} (ghaṭṇā)}{ਘਟਣਾ (ghaṭṇā)}]{\pa{ਘਟਣਾ} (ghaṭṇā) — to go down; to decrease}

\label{lesson:PA-C290-ghatna}



\begin{quote}
Before the new one: say the Punjabi for hail, then the Punjabi for to clear up.
\end{quote}

\subsection*{You'll want to know: \pa{ਘਟਣਾ}}
\textbf{\pa{ਘਟਣਾ}} — \emph{ghaṭṇā} — \textquotedblleft{}to go down; to decrease\textquotedblright{}.

\begin{quote}
\pa{ਸ਼ਾਮ} \pa{ਨੂੰ} \pa{ਗਰਮੀ} \pa{ਘਟੇਗੀ।}
\end{quote}

\begin{quote}
\emph{shām nū̃ garmī ghaṭegī.}
\end{quote}

\begin{quote}
The heat will go down in the evening.
\end{quote}

Past (it …): \pa{ਉਹ} \pa{ਪਾਣੀ} \pa{ਘਟਿਆ}, \emph{oh pāṇī ghaṭiā}.

\begin{cousinweb}
\textbf{\pa{ਘਟਣਾ}} (\emph{ghaṭṇā}) is inherited from a reconstructed \emph{*ghaṭṭati}, \textquotedblleft{}decreases, is wanting\textquotedblright{}; CDIAL lists Punjabi \emph{ghaṭṇā}, \textquotedblleft{}to decrease\textquotedblright{}, under it (CDIAL 4415).
\end{cousinweb}

\begin{grammarlens}[title={What you've built}]
One more word for saying what comes next.
\end{grammarlens}

\subsection*{Your turn}
\begin{itemize}
\raggedright
  \item \emph{Say it:} \emph{ghaṭṇā}
  \item \emph{Say it:} \emph{ghaṭṇā}, once more
  \item \emph{Say it:} \emph{shām nū̃ garmī ghaṭegī.} — The heat will go down in the evening.
  \item \emph{Recall:} say the Punjabi for hail, then the Punjabi for to clear up, then say \emph{ghaṭṇā} again
  \item \emph{Write it:} copy \textbf{\pa{ਘਟਣਾ}} once, checking the whole word
\end{itemize}

\begin{tcolorbox}[breakable,colback=teal!4,colframe=teal!35!black,title={Before you move on}]
What does \pa{ਘਟਣਾ} mean? (\textquotedblleft{}To go down; to decrease\textquotedblright{}.) Say it once more.
\end{tcolorbox}
\section[\texorpdfstring{\pa{ਵਗਣਾ} (vagṇā)}{ਵਗਣਾ (vagṇā)}]{\pa{ਵਗਣਾ} (vagṇā) — to blow; to flow}

\label{lesson:PA-C290-vagna}



\begin{quote}
Before the new one: say the Punjabi for to clear up, then the Punjabi for to go down.
\end{quote}

\subsection*{You'll want to know: \pa{ਵਗਣਾ}}
\textbf{\pa{ਵਗਣਾ}} — \emph{vagṇā} — \textquotedblleft{}to blow; to flow\textquotedblright{}.

\begin{quote}
\pa{ਰਾਤ} \pa{ਨੂੰ} \pa{ਤੇਜ਼} \pa{ਹਵਾ} \pa{ਵਗੇਗੀ।}
\end{quote}

\begin{quote}
\emph{rāt nū̃ tez havā vagegī.}
\end{quote}

\begin{quote}
A strong wind will blow at night.
\end{quote}

Past (it …): \pa{ਉਹ} \pa{ਹਵਾ} \pa{ਵਗੀ}, \emph{oh havā vagī}.

\begin{cousinweb}
\textbf{\pa{ਵਗਣਾ}} (\emph{vagṇā}) is inherited from Sanskrit \emph{vrajati}, \textquotedblleft{}goes, wanders\textquotedblright{}; CDIAL lists Punjabi \emph{vagṇā}, \textquotedblleft{}to move, flow, blow\textquotedblright{}, under it (CDIAL 12225).
\end{cousinweb}

\begin{grammarlens}[title={What you've built}]
One more word for saying what comes next.
\end{grammarlens}

\subsection*{Your turn}
\begin{itemize}
\raggedright
  \item \emph{Say it:} \emph{vagṇā}
  \item \emph{Say it:} \emph{vagṇā}, once more
  \item \emph{Say it:} \emph{rāt nū̃ tez havā vagegī.} — A strong wind will blow at night.
  \item \emph{Recall:} say the Punjabi for to clear up, then the Punjabi for to go down, then say \emph{vagṇā} again
  \item \emph{Write it:} copy \textbf{\pa{ਵਗਣਾ}} once, checking the whole word
\end{itemize}

\begin{tcolorbox}[breakable,colback=teal!4,colframe=teal!35!black,title={Before you move on}]
What does \pa{ਵਗਣਾ} mean? (\textquotedblleft{}To blow; to flow\textquotedblright{}.) Say it once more.
\end{tcolorbox}
